% Nota de comparacao_criterios. GERADA por flim_al/tabelas.py.
\footnotesize \texttt{random} é o piso --- sem critério nenhum. O pareamento casa split, decoder, orçamento e semente \textbf{dentro da mesma família de experimento}; a média mostrada vem dessas mesmas células, para que média e $\Delta$ falem do mesmo conjunto. \textbf{Toda técnica que rodou está aqui, inclusive as que perderam do sorteio} --- no experimento de encoder retreinado isso aconteceu com as quatro de imagem. Negrito só com p < 0,05. Para técnicas de REGIÃO, leia a coluna do controle de nível (\texttt{random\_region}): o $\Delta$ contra \texttt{random} soma seleção e geometria, e a geometria pesa ~2$\times$ a seleção neste projeto. \textbf{\texttt{flim\_3img} é o FLIM do artigo} --- as 3 imagens que os autores escolheram, sem seleção automática. A coluna \texttt{$\Delta$ vs FLIM} responde à pergunta central: aplicar AL melhora sobre isso? Ela só aparece em orçamento casado (K=3), porque o braço do artigo só existe com 3 imagens; comparar AL com 10 imagens contra FLIM com 3 mediria orçamento, não seleção. \textbf{Coluna toda vazia significa comparação IMPOSSÍVEL com os dados existentes, não ausente por descuido}: o braço do artigo usa os 31 markers REAIS do especialista, e os braços de AL sobre o pool completo usam sintéticos, porque marker real só existe para 31 imagens. \texttt{marker\_origem} entra no pareamento para que esse par não se forme --- ele mediria quem desenhou o traço e apresentaria como efeito da seleção. Para responder à pergunta, rode \texttt{scripts/varrer\_criterios.py --modo imagem}: marker real nos dois lados, com \texttt{oracle} (o passo 7 do Algoritmo 1) como referência.
\normalsize
